% Workshop handout with a restrained cover, code sample and workflow diagram.
% Tectonic reads UTF-8 directly. Replace each sample passage and comment.
\documentclass[12pt]{article}

\usepackage[{{language}}]{babel}
\usepackage[T1]{fontenc}
\usepackage{newpxtext}
\usepackage{microtype}
\usepackage[skip=0.55\baselineskip]{parskip}
\usepackage[a4paper,margin=2.5cm,headheight=14pt,headsep=12pt]{geometry}
\usepackage{xcolor}
\usepackage{enumitem}
\usepackage{fancyhdr}
\usepackage{titlesec}
\usepackage{tcolorbox}
\usepackage{booktabs}
\usepackage{graphicx}
\usepackage{float}
\usepackage{hyperref}

\definecolor{workshopink}{HTML}{292927}
\definecolor{workshopaccent}{HTML}{8A493D}
\hypersetup{colorlinks=true,linkcolor=workshopaccent,urlcolor=workshopaccent,
  citecolor=workshopaccent,breaklinks=true,pdftitle={{{title}}},pdfauthor={{{author}}}}

% Section numbers and rules use a single warm accent; body copy stays dark.
\titleformat{\section}{\normalfont\Large\bfseries\color{workshopink}}
{\textcolor{workshopaccent}{\thesection}}{0.75em}{}
\titlespacing*{\section}{0pt}{1.4\baselineskip}{0.45\baselineskip}
\titleformat{\subsection}{\normalfont\large\bfseries\color{workshopink}}
{\thesubsection}{0.7em}{}
\setlength{\parindent}{0pt}
\setlength{\emergencystretch}{2em}
\setlist{leftmargin=1.4em,itemsep=2pt,topsep=4pt}

% Running head: deliverable at left, course at right, page number below.
\pagestyle{fancy}
\fancyhf{}
\lhead{{{deliverable}}}
\rhead{{{course}}}
\cfoot{\thepage}
\renewcommand{\headrulewidth}{0.5pt}
\renewcommand{\headrule}{\hbox to\headwidth{\color{workshopaccent}\leaders\hrule height \headrulewidth\hfill}}

\begin{document}

  % Cover page: keep this block to one page.
  \begin{titlepage}
    \centering
    {\scshape\large {{deliverable}} \par}
    \vspace{1.1cm}
    {\huge\bfseries\color{workshopink} {{title}} \par}
    \vspace{0.45cm}
    {\Large {{subtitle}} \par}
    \vspace{1.5cm}
    \begin{tcolorbox}[colback=workshopaccent!4,colframe=workshopaccent!55,
        width=0.88\textwidth,boxrule=0.6pt,arc=1mm]
      \begin{tabular}{@{}ll@{}}
        \textbf{Student:} & {{author}} \\
        \textbf{Course:} & {{course}} \\
        \textbf{Deliverable:} & {{deliverable}} \\
        \textbf{Date:} & {{date}} \\
        \textbf{Repository:} & {\small\url{{{repository}}}} \\
      \end{tabular}
    \end{tcolorbox}
    \vfill
  \end{titlepage}

  \section{Problem Statement}
  % Replace with the official prompt and expected outcome.
  The task is to design and implement an algorithm that solves a concrete
  problem with conditional and iterative steps. Describe the approach, write
  the code and check it against representative cases. Project updates are
  available at \url{{{repo_url}}}.

  \section{Problem Context}
  % Replace with the context and constraints that motivate your solution.
  The current process is manual and takes time: each case must be checked
  separately, and the results are not always consistent. Automating the
  process should make the outcome repeatable and easier to review.

  \section{Algorithmic Solution}
  % Replace with your inputs, outputs and key steps.
  The solution reads an input, validates it, applies a simple rule and
  presents the result in an organised form.

  \begin{tcolorbox}[colback=gray!4,colframe=gray!45,coltitle=workshopink,
      title={Algorithm Steps},boxrule=0.5pt,arc=1mm]
    \begin{enumerate}
      \item Initialise the data and state.
      \item Read the user's input.
      \item Validate the input and detect errors.
      \item Process the data according to the problem rules.
      \item Present the results clearly.
    \end{enumerate}
  \end{tcolorbox}

  % Replace the nodes with your own flow; keep the style and caption.
  \begin{mermaid}[style=editorial,width=0.68\linewidth,
      caption={Algorithm workflow}]
    flowchart TD
    A[Start] --> B[Read input]
    B --> C{Input valid?}
    C -- No --> D[Show error]
    C -- Yes --> E[Process data]
    E --> F[Show result]
    D --> G[Finish]
    F --> G
  \end{mermaid}

  \section{Implementation}
  % Replace with the core of your solution (up to 20 lines).
  The listing shows the core transformation from input values to results.

  \begin{code}[lang=python,numbers=true]
def classify(values):
    even, odd = [], []
    for value in values:
        (even if value % 2 == 0 else odd).append(value)
    return even, odd

def parse_values(text):
    if not text.strip():
        raise ValueError("input is empty")
    return [int(value) for value in text.split(",")]
  \end{code}

  \section{Tests}
  % Replace the examples with the cases you actually ran.
  The table records representative inputs, expected outcomes and status.

  \begin{table}[H]
    \centering
    \caption{Representative algorithm test cases}
    \begin{tabular}{@{}llll@{}}
      \toprule
      Case & Input & Expected output & Status \\
      \midrule
      Typical & 8 & Even: 8 & Pass \\
      Boundary & Empty list & No values & Pass \\
      Invalid & Non-numeric text & Error message & Pass \\
      \bottomrule
    \end{tabular}
  \end{table}

  \section{Conclusion}
  % Replace with what worked and what remains to be done.
  The algorithm solves the stated task and passes the selected test cases.
  Its steps remain separate, so adding a new rule does not require rewriting
  the entire implementation.

\end{document}
